\begin{papercard}{messner2020}
\cf{Stages} \stagetag{A} \stagetag{W} \stagetag{I}. Translation and validation of the German Mobile App Rating Scale (MARS-G). A model for instrument translation, rater training and reliability reporting in the German setting.

\cf{Aim} Translate and validate a German MARS and test the scalability of the overall score with Mokken scale analysis \pc{messner2020}{2}{aimed to develop and validate a German version of the Mobile}.

\cf{Translation procedure} Two independent forward translations by bilingual scientists, then discussion to a pilot version. Blind back-translation by two bilingual people of different backgrounds (a postdoctoral psychologist and a non-academic). Comparison with the original produced a penultimate version, which was checked for comprehensibility by three researchers and three non-academics before finalisation \pc{messner2020}{2}{The MARS was translated from English into German by two}.

\cf{Sample and search} Validated inside a review of anxiety apps. Thirteen search terms were entered \emph{separately}, because the stores allow no truncation or Boolean operators \pc{messner2020}{3}{Each search term was provided}. Three steps: search, screening of store pages, and eligibility after download (anxiety focus, German or English, official store, no app bundles, no device incompatibility or test phase) \pc{messner2020}{3}{The inclusion process was divided into three steps}. 3{,}562 hits, minus 810 duplicates, 2{,}577 inappropriate on screening and 71 ineligible, left 104 apps \pc{messner2020}{3}{We identified 3562 MHAs from the app stores}.

\cf{Rater training} Followed Stoyanov's original procedure. A training video demonstrated rating on an example app, and each rater then rated five predefined apps. Deviations of $\geq$2 points from the standard rating were discussed to agreement. Raters held at least a bachelor's degree in psychology \pc{messner2020}{3}{Each rater was trained using this}. Each app was used for 15 minutes and rated immediately after, in both language versions \pc{messner2020}{3}{The raters tested all MHAs for 15 minutes}.

\cf{Instrument additions} The description section gained: theoretical background, methods (incl.\ tracking and feedback), store category, embedding in routine care, type of use, guidance, certification (medical device law) and data safety (log in, informed consent). A new section rates therapeutic gain: gain for patient and therapist, risks and adverse effects, ease of implementation \pc{messner2020}{3}{We added the following items in the app description section for}. The German form's data-protection block is a yes/no list: password, login, privacy policy (\emph{Datenschutzerklärung}), active consent, conflicts of interest, contact/imprint, secure transmission, storage location, sharing with third parties, emergency functions, and strategies for phone loss \pc{messner2020marsg}{3}{Sicherheit und Datenschutz}. Quality items are rated 1--5 with anchored descriptions \pc{messner2020marsg}{4}{Das MARS Rating erfasst die Qualit}.

\cf{Psychometrics} ICC bands: $<$0.50 poor, 0.51--0.75 moderate, 0.76--0.89 good, $>$0.90 excellent; $>$0.75 counted as sufficient \pc{messner2020}{3}{of <0.50 indicated poor correlation, 0.51-0.75 indicated}. Internal consistency was measured with McDonald's $\omega$ rather than Cronbach's $\alpha$, using robust, bootstrapped CIs \pc{messner2020}{3}{Omega was used instead of the widely adopted}. Validity: MARS vs.\ MARS-G correlations with an \emph{a priori} threshold $r>0.8$, plus $t$-tests \pc{messner2020}{3}{value >0.8 was a priori considered by the author group as an}. Mokken scale analysis (Loevinger's $H$) \pc{messner2020}{3}{Mokken scale analysis (MSA) is a scaling approach closely}.

\cf{Results} ICC 0.84 (MARS) and 0.83 (MARS-G) \pc{messner2020}{4}{The ICCs for the MARS and MARS-G were high}. Dimension correlations 0.93--0.98. No association between store star ratings and quality \pc{messner2020}{5}{were no associations between user ratings and quality ratings}. Planned missingness from ``not applicable'' (item 19 applied to only one app) distorted the $\omega$ confidence interval for information quality \pc{messner2020}{5}{the confidence interval of omega was}. Recommendation: train with the five-app set and revise ratings until ICC $>$0.75 \pc{messner2020}{5}{The MARS-G ratings should be revised until an}.

\cf{Limitations and future work} Convergent validity was tested only against the MARS. Anxiety apps only. No CFA \pc{messner2020}{6}{First, convergent validity}. A key methods observation: store result order is ``incomprehensible'' and differs depending on the account used, so they suggest a web crawler over European stores \pc{messner2020}{6}{the sequence in which apps are presented}.

\cf{Use for this thesis} (1) A German-language precedent for structured app appraisal and its privacy checklist wording (\emph{Datenschutzerklärung}, \emph{Impressum}, Ort der Datenspeicherung). (2) A complete template for translating any English instrument item: forward, back and comprehensibility check. (3) A rater training protocol with calibration apps and a deviation rule. (4) Evidence that the store ranking depends on the account, which supports a logged-out/fresh-account search and recording the account state. MARS measures app quality, not harm communication.
\end{papercard}


\begin{papercard}{messner2020marsg}
\cf{What it is} The 15-page German rating form (``MARS-German'') attached to Messner et al.\ (2020), a translation of the MARS with added sections for psychotherapy apps \pc{messner2020marsg}{1}{Mobile Anwendungen Rating Skala}. It is a form, not a study: the card lists what the form asks and where it differs from the English MARS. The classification section collects descriptive and technical data from the store page \pc{messner2020marsg}{1}{Diese Sektion dient dazu, beschreibende und technische Informationen}; fields cover name, ratings of this and all versions with their counts, developer, version, last update, costs of basic and extended version, platform, affiliation (unknown, commercial, government, NGO, university), focus, theoretical background, methods (23 options), age group, technical aspects, store category, embedding into therapy, field of use, type of treatment, level of support and certification.

\cf{Security and privacy checklist} A yes/no block on page 3 asks: password use allowed, login required, privacy policy present, active confirmation of a consent form required \pc{messner2020marsg}{3}{Erfordert das aktive Best}, funding or conflict-of-interest notes, contact or imprint, security of data transmission \pc{messner2020marsg}{3}{Sicherheit der Daten}, place of data storage \pc{messner2020marsg}{3}{Ort der Datenspeicherung}, data shared with third parties \pc{messner2020marsg}{3}{Werden Daten mit Dritten geteilt}, emergency functions, security strategies when the phone is lost \pc{messner2020marsg}{3}{Sicherheitsstrategien bei}, and other. A certification field offers ``none'', ``conforms to the Medical Devices Act'' or other \pc{messner2020marsg}{3}{Entspricht dem Medizinproduktegesetz}. All of this is presence/absence, with no anchors for quality and no link to the policy text.

\cf{Rating sections} Sections A--D (engagement, functionality, aesthetics, information) use anchored 1--5 descriptors per item with a section mean that excludes ``not available'' answers; section E (items XXIV--XXVII) is subjective quality (recommendation, expected use, willingness to pay, stars); section F gives six items on perceived impact on awareness, knowledge, attitude, intention, help seeking and behaviour. The added section PT (psychotherapy) has four items (XX--XXIII) on patient benefit, therapist benefit (blended care only), possible risks and harmful effects, and transferability to routine care \pc{messner2020marsg}{12}{Diese Sektion erfasst G}. The risk item has five anchored levels from ``harmful from a therapeutic view, high risk of undesired effects (e.g.\ misleading advice, harmful exercises, omitted help)'' to ``safe and in line with current best-practice guidelines'' \pc{messner2020marsg}{12}{Die App ist aus therapeutischer Sicht sch}, and a footnote warns strongly against using apps scoring below 3 \pc{messner2020marsg}{12}{Von der Benutzung von Apps die hier einen Wert unter 3 erreichen wird dringend abgeraten}. Level 3 gives ``no emergency function'' as an example of a shortfall against best practice. Transferability is anchored by whether the app was tested on patients and under representative routine conditions \pc{messner2020marsg}{13}{in die Routineversorgung}.

\cf{Limits (card)} The form's harm item is clinical (treatment risk), not data-protection harm; privacy appears only in the yes/no checklist; the German wording of anchors cannot be quoted for the English thesis without translation notes; the form carries no reliability data itself.

\cf{Use for this thesis} (1) German-language store fields (store category, certification, embedding in therapy) for the app sheet in the German setting, including the DiGA context. (2) A model of an \emph{anchored risk item} with a stop rule (below 3: do not use), which can inspire an anchored harm-severity or harm-communication scale; the thesis would need privacy-specific anchors. (3) The checklist items on transmission security, storage location and third-party sharing show the minimum data-protection facts that a rater could verify from a store page and policy.
\end{papercard}
